Um experimento para determinar o período natural ou fundamental de oscilação de um sistema simples de massa-mola (ver figura P3.34) é estabelecido da seguinte forma. Uma mola de rigidez $k$ é fixada em uma extremidade e conectada a uma massa $m$ na outra, com a massa sendo capaz de se mover ao longo de uma pista de ar ideal e sem atrito. A massa é deslocada a uma distância $x_0$ de sua posição inicial de repouso, após o que oscila ao longo da pista de ar em torno dessa posição inicial. O tempo necessário para uma oscilação completa, isto é, o período, é medido várias vezes por vários períodos sucessivos, com os resultados sendo comparados com a fórmula teórica para o período, $T$: $T = 2\pi\sqrt{m/k}$. Supondo que $k$ e $m$ sejam conhecidos e que o cronômetro usado para medir o período tenha precisão de $\pm 1\%$. Quais são as possíveis armadilhas que podem impedir a determinação experimental bem-sucedida de $T$?

